\section{Fixed representations and forecasting procedures}
\label{sec:representations}
The history length, forecast horizon, observed target, and origin rule define the task before representations are compared. Development comprises feature fitting, head fitting, and permitted model selection. Population calibration uses independent units after the candidate family is frozen. Chronological assessment on a single observed series has a different evidential role and is reported as a finite-period comparison.

\subsection{KAHM as an expandable example}
KAHMs use local kernel reconstruction models fitted on training histories \citep{kumar2024kahm,kumar2025collaborative,kumar2026limits}. For a preprocessed history $z(x)$ and reconstruction $\widehat z_{bcj}(x)$ in reference group $c$ of retained feature block $b$, write
\[
e_{bcj}(x)=1-\exp\{-\|z(x)-\widehat z_{bcj}(x)\|_2\},
\qquad
\theta_{bcj}(x)=\pi^{-1}\arccos\gamma_{bcj}(x).
\]
Here $\gamma$ is clipped cosine similarity, with the declared zero-norm convention. The group discrepancy and normalized association are
\begin{align*}
T_{bc}(x)&=\min_j\sqrt{\{e_{bcj}(x)^2+\theta_{bcj}(x)^2\}/2},\\
\Phi_{bc}(x)&=
\frac{(1-T_{bc}(x)+\tau)^\omega}
{\sum_r(1-T_{br}(x)+\tau)^\omega},
\qquad \tau>0,\ \omega>0.
\end{align*}
These simplex coordinates are reconstruction affinities, not calibrated response probabilities. Stable log-space normalization is part of the implementation. The local PCA request is
$r_0=\min(r_{\max},d_x,N_{\rm loc}-1)$, followed by the inherited variance and projected-range rules; requested cap and stored rank are different quantities.

Appending fixed coordinate blocks produces $\Gamma_s=(\Phi_1,\ldots,\Phi_s)$. Under the maximum of within-block Euclidean distances, the product of simplices is compact with diameter at most $\sqrt2$, independently of the number of blocks. Thus the scalar hierarchy applies to any bounded scalar target under unconditional iid sampling, even if the history-to-feature map is discontinuous. Adding feature blocks changes the candidate class; increasing the number of observations in a certificate block changes the relaxation for a fixed class.

Training-only min--max and standard scaling are both retained where declared. Scaling can change the radial and angular terms, learned partitions, and local ranks. It must not be treated as an innocuous coordinate change while keeping the same numerical sensitivity budget. For example, replacing $U$ by $cU$ preserves a Lipschitz class only after the metric or budget is compensated. Likewise, for fixed positive scores, changing the concentration $\omega$ is an invertible reparameterization in exact arithmetic, yet can alter accessibility to a restricted decoder and numerical conditioning.

\subsection{Theoretical family and implemented heads}
A retained-block affine head with $\sum_b\|W_b\|_{\rm F}\leq\kappa$ is $\kappa$-Lipschitz in the maximum-block metric. Projection onto a known output box is nonexpansive. Zero padding retains earlier feasible forecasts, while merging coefficients of identical blocks proves invariance to exact block duplication. These are statements about feasible classes. A ridge penalty can change under duplication, and a fitted predictor need not improve when its class grows.

The independent-episode controlled study uses this retained-block family. The later chronological studies instead declare fixed Euclidean feature metrics and their own coefficient constraints. Their finite-cohort results are not silently transferred to the retained-block population theorem. The final affine protocol uses a rational head, a Euclidean norm bound of two, exact intercept recentering, and normalized output clipping. Its training-derived output range defines that predictor; it is not assumed to contain every future target.

Costs must also be specified. Coordinate count can order an abstract family, but it does not establish an ordering of complete storage or runtime. When deployment utility is assessed, feature extraction, stored encoder state, head evaluation, checking, and any cached state must be included or explicitly separated.

